\section*{Introduction}

The use of pointclouds has steadily increased over the years, driven by the growing availability of 3D acquisition systems and advanced 3D processing and visualization software. Pointclouds are nowadays used in many areas, including engineering and manufacturing, architecture and construction, cultural heritage documentation and preservation, robotics and autonomous systems, extended and immersive reality, geospatial and environmental analysis, to name just a few.

These tutorials are for those wishing to learn a little bit more about the basics of pointcloud processing. Having gone through this stage during my Ph.D., I hope here to share some of what I have learned so far.

The notebooks are designed to make pointcloud processing algorithms easier to understand, without compromising performance too much and trying to minimize the use of specialized third-party software. They require basic knowledge of Python and its main scientific libraries.

\subsection*{Dependencies}

The code is in python and relies on \texttt{numpy}, \texttt{scipy}, \texttt{matplotlib}, \texttt{scikit-learn}, and \texttt{jupyterlab}.

These dependencies may be installed with pip (or conda) with

\texttt{pip install numpy scipy matplotlib scikit-learn jupyterlab}

JupyterLab may be started using the terminal or Anaconda prompt simply by typing 

\texttt{jupyter lab}

\subsection*{Reusing and distributing}

This notebook is licensed under the Creative Commons Attribution-NonCommercial 4.0 International License (CC-BY-NC-4.0). For the full license text, please visit \url{https://creativecommons.org/licenses/by-nc/4.0/}.

If you use these tutorials in your work, please cite them as:

\begin{Shaded}
\begin{Highlighting}[]

@unpublished\{gregorio2024tutorials,
     author=\{Grégorio, Jean-Loup\},
     title=\{Tutorials for pointcloud processing in Python\},
     year=\{2024\},
\}

\end{Highlighting}
\end{Shaded}
